\documentclass[10pt]{article}
\usepackage[margin=.9in]{geometry}
\usepackage[T1]{fontenc}
\usepackage{lmodern,amsmath,amssymb,amsthm,mathtools,booktabs,microtype}
\usepackage[colorlinks=true,linkcolor=blue]{hyperref}
\newtheorem{theorem}{Theorem}[section]
\newtheorem{lemma}[theorem]{Lemma}
\title{Canonical ReLU Swing-by: certified conditional timing}
\author{}\date{}
\begin{document}\maketitle
\noindent
The full-tube conditional timing theorem is closed. It separates the
repository statistic's sign transition from every full OOD-error minimizer
on the certified time window, uniformly over an explicit sector. Exact
initialization-to-entry reachability remains open. The exact fixed-pattern
moment equations and their event-transfer limitation are stated separately.
% Conditional on the existing initialized-certificate entry, not on timing signs.
\section{Conditional canonical timing and minimizer localization}
\label{sec:canonical-full-tube-timing}
Use the exact binary canonical dataset, exact archived reference, and fixed
71-neuron ID-selected cohort of the existing continuation certificate.
The loss is $\mathcal L=(2N)^{-1}\sum_n\|f_\theta(x_n)-x_n\|^2$,
with $N=4000$, $h=200$. All times below are physical times. Set
\[
 r=\sin(\pi/36),\quad k=\cos(\pi/36),\quad
 \xi_\alpha=(\cos\alpha,\sin\alpha),\quad
 \mathcal I=\{\alpha:|\alpha+17\pi/36|<10^{-6}\}.
\]
For the fixed ungated strong block $M_I=(m_{ab})$, retain exactly the
repository statistic
\[
 Q_G=[(m_{11}-1)r-m_{12}](r\dot m_{11}-\dot m_{12})
           +r\dot m_{21}-\dot m_{22}.
\]
The factors $k$ are not inserted into this auxiliary statistic.
Write $q_E(\alpha,t)=\frac{d}{dt}\frac12\|f_t(\xi_\alpha)-\xi_\alpha\|^2$,
with time derivatives understood almost everywhere.

\begin{theorem}[Conditional sign separation and minimizer localization]
\label{thm:canonical-full-tube-timing}
Every global Clarke flow satisfying the existing entry condition
\[
 \|\theta(3.4)-\widehat\theta_{3.4}\|_2\le1.2\cdot10^{-5}
\]
has the following properties. Put
\[
 [g_-,g_+]=[3.5738,3.6348],\qquad
 [e_-,e_+]=[3.6384,3.7152].
\]
For every compatible training-selector realization,
\begin{align*}
 Q_G&<0&&\text{a.e. on }[3.4,g_-],&
 Q_G&>0&&\text{a.e. on }[g_+,3.9],\\
 q_E(\alpha,\cdot)&<0&&\text{a.e. on }[3.4,e_-],&
 q_E(\alpha,\cdot)&>0&&\text{a.e. on }[e_+,3.9]
 \qquad(\alpha\in\mathcal I).
\end{align*}
Consequently every minimizer of $E(\xi_\alpha,\cdot)$ on $[3.4,3.9]$
lies in $[e_-,e_+]$, uniformly for $\alpha\in\mathcal I$.
The sign-transition brackets are strictly separated:
\[
 e_- -g_+=\frac{18}{5000}=0.0036>0.
\]
Throughout the interval $[g_+,e_-]$, almost everywhere and uniformly in
$\alpha\in\mathcal I$,
\[
 Q_G>1.21\cdot10^{-7},\qquad
 q_E(\alpha,\cdot)<-5.44\cdot10^{-7},\qquad
 q_E(\alpha,\cdot)-Q_G<-2.047\cdot10^{-4}.
\]
For the central probe $\alpha=-17\pi/36$, the tighter minimizer bracket is
$[3.6394,3.7136]$, giving separation $23/5000=0.0046$.
No uniqueness or continuity of a $Q_G$ zero, unique OOD minimizer, or
initialized reachability is asserted.
\end{theorem}

\begin{proof}
The existing flow certificate is used without modification. On each of its
2500 intervals $[t_j,t_{j+1}]$, its nonnegative scalar comparison is
nondecreasing, so its terminal bound $r_{j+1}$ bounds the distance to the
cubic reference throughout that entire interval. This justifies using the
recorded \texttt{radius\_upper}, rather than the larger prescribed guard
radius. No flow calculation is rerun.

Let $y_j$ be the exact binary midpoint used by the timing verifier. If
$d_j$ bounds motion of the cubic reference about its interval midpoint,
and $\epsilon_j$ bounds midpoint rounding, then
\[
 \|\theta(t)-y_j\|\le R_j:=r_{j+1}+d_j+\epsilon_j
 \qquad(t_j\le t\le t_{j+1}).
\]
The cubic derivative bounds and all these additions are outward-rounded.
For $P_j=\|y_j\|$, $\bar P_j=P_j+R_j$, the static-neighborhood proposition
in the continuation certificate provides a polynomial-field Lipschitz bound
$L_j$ and an actual-training-gate correction $B_j$. The latter encloses
all compatible selectors, including every ambiguous training gate.

For either potential, put $g=\nabla\Phi$, $H\ge\sup\|Dg\|$ on this ball.
The checked static inequality is
\[
 |g(\theta)^TF(\theta)-g(y_j)^TF(y_j)|
 \le R_j\{H(\|F(y_j)\|+L_jR_j)+\|g(y_j)\|L_j\}
       +(\|g(y_j)\|+HR_j)B_j.
\]
For $Q_G$, $\Phi=\tfrac12[((m_{11}-1)r-m_{12})]^2+rm_{21}-m_{22}$.
For the full error, use the actual probe-active matrix and the quadratic
reconstruction potential. For the discrepancy, apply the same inequality
\emph{directly} to $g_E-g_G$, with $H_E+H_G$ as a valid derivative bound;
separately enlarged rate intervals are not subtracted.

Probe signs are certified on each entire tube. For each reference
preactivation $q$ its four Bernstein control coefficients are
\[
 q_0,\quad q_0+\tfrac{1}{3}h\dot q_0,\quad
 q_1-\tfrac{1}{3}h\dot q_1,\quad q_1,\qquad h=1/5000.
\]
Their interval convex hull bounds the exact reference cubic. Subtract the
parameter-tube perturbation $r_{j+1}\|\xi\|$ and the sector cost
$(\sup_t\|\widehat u_i(t)\|+r_{j+1})10^{-6}$ from the oriented margin.
The minimum resulting margin over all steps and neurons exceeds
$1.40926\cdot10^{-5}$. Thus no probe gate is unresolved. The smooth probe
potential may be extended polynomially to the larger midpoint ball even
when only the reference tube is known to preserve its gates: equality with
the actual error is required only on the actual tube.

On each tube the angular transport cost for the error rate is bounded by
\[
 2(1+\bar P_j^2/2)\bar P_j\sup\|F\|\,10^{-6}.
\]
Indeed the probe-cell matrix has norm at most $\bar P_j^2/2$, its time
derivative at most $\bar P_j\|F\|$, and differentiation of
$\langle M\xi-\xi,\dot M\xi\rangle$ in $\xi$ gives this Lipschitz bound.
$Q_G$ is independent of the varied probe angle, so the same cost applies
to $q_E-Q_G$.

All 2500 interval enclosures pass. The negative-prefix and positive-suffix
margins, rounded outward, are respectively
\begin{center}
\begin{tabular}{lcc}
\toprule
Quantity & negative prefix: $-\text{upper}$ & positive suffix: lower\\
\midrule
$Q_G$ & $>1.0839\cdot10^{-7}$ & $>1.2158\cdot10^{-7}$\\
$q_E$, central probe & $>3.2526\cdot10^{-7}$ & $>3.5073\cdot10^{-7}$\\
$q_E$, entire sector & $>5.4476\cdot10^{-7}$ & $>1.9957\cdot10^{-8}$\\
\bottomrule
\end{tabular}
\end{center}
The sign-unresolved steps are enclosed in exactly the stated brackets.
The 18 whole steps between $g_+$ and $e_-$ have both required strict signs.
On these steps the direct central discrepancy enclosure is contained in
$[-3.05302\cdot10^{-4},-2.08690\cdot10^{-4}]$; after angular transport its
upper bound is less than $-2.04714\cdot10^{-4}$.
The exact step ledger, every sign-unresolved index, input hashes, and
neighboring strict margins appear in
\texttt{TIMING\_FULL\_TUBE\_CERTIFICATE.csv} and its JSON summary.

The error is absolutely continuous and hence continuous on the compact
interval. Integrating the strictly negative prefix bound excludes any
minimizer before $e_-$, and integrating the positive suffix bound excludes
any minimizer after $e_+$. The $Q_G$ sign bounds establish an earlier
negative-to-positive sign transition bracket, without invoking the
intermediate value theorem for a discontinuous training field.
\end{proof}

\paragraph{Independent interior witness.}
The proposed $t=3.64$ witness is independently reproduced. In the static
ball of radius $2.6457445093\cdot10^{-5}$ (enclosed upward), outward-rounded
bounds are
\begin{align*}
 Q_G&\in[2.7405\cdot10^{-5},2.5834\cdot10^{-4}],\\
 q_E&\in[-2.1933\cdot10^{-4},-8.7767\cdot10^{-6}],\\
 q_E-Q_G&\in[-3.0089\cdot10^{-4},-2.1296\cdot10^{-4}].
\end{align*}
The speed is below $0.330122389$. The certified radius at $3.64$ leaves
more than $2\cdot10^{-6}$ available, while $6\cdot10^{-6}$ time units cost
less than $1.980735\cdot10^{-6}$ in parameter distance. A first-exit
argument gives $Q_G>0$ and $q_E<0$ on $[3.64,3.640006]$ for the central
probe. The displayed rounded intervals are deliberately outward.

\paragraph{Scope and provenance.}
The theorem concerns the supplied exact binary dataset and every entering
Clarke flow; the archived exact initial weights have not been connected to
this entry set. The historical archive has no independent raw-data hash.
The supplied binary data match the pinned repository generator under the
recorded runtime; this is not a high-probability theorem. Neither the
empirically interpolated times nor a uniqueness assumption is used.

\section{Exact dynamics of training-pattern moments}
\label{sec:pattern-quotient}
The residual convention matters. In this section define
$\rho_n=f_\theta(x_n)-x_n$ and
$A_g=N^{-1}\sum_n g_n\rho_nx_n^T$.
If instead the residual is $x_n-f_\theta(x_n)$, all three displayed minus
signs below become plus signs. Both conventions describe the same flow.

\begin{lemma}[Fixed-pattern moment dynamics]\label{lem:pattern-quotient}
On an open time interval where every training preactivation is nonzero
and the neuronwise training patterns are fixed, let
\[
 I_g=\{i:\mathbf1_{u_i^Tx_n>0}=g_n\text{ for all }n\},\quad
 M_g=\sum_{i\in I_g}w_iu_i^T,\quad
 U_g=\sum_{i\in I_g}u_iu_i^T,\quad
 W_g=\sum_{i\in I_g}w_iw_i^T.
\]
Then the exact closed equations are
\begin{align*}
 \dot M_g&=-A_gU_g-W_gA_g,\\
 \dot U_g&=-A_g^TM_g-M_g^TA_g,\\
 \dot W_g&=-A_gM_g^T-M_gA_g^T,
 \qquad f(x_n)=\sum_g g_nM_gx_n.
\end{align*}
All moments include their entire groups; no concentration or low-rank
approximation is assumed.
\end{lemma}
\begin{proof}
Direct differentiation of the stated normalized loss gives
$\dot w_i=-A_gu_i$ and $\dot u_i=-A_g^Tw_i$ for $i\in I_g$.
Differentiate each outer product and sum. For example
$\dot w_iu_i^T+w_i\dot u_i^T=-A_gu_iu_i^T-w_iw_i^TA_g$.
The other two identities follow in the same way, with the transposes shown.
The network identity follows from $(u_i^Tx_n)_+=g_nu_i^Tx_n$.
Consequently $A_g$ depends only on the fixed patterns, the data, and the
moments, proving closure on this interval.
\end{proof}

Put
\[
 K_g=\begin{pmatrix}U_g&M_g^T\\M_g&W_g\end{pmatrix},\qquad
 B_g=\begin{pmatrix}0&A_g^T\\A_g&0\end{pmatrix}.
\]
Then $K_g\succeq0$, $\dot K_g=-B_gK_g-K_gB_g$ and
$\operatorname{tr}(U_g-W_g)$ is constant on the fixed-pattern interval.
The same equations hold for a partition refined by fixed cohort membership
and fixed probe signs, since every subgroup with training pattern $g$
uses the same $A_g$.

\paragraph{What the quotient retains and loses.}
The training-pattern moments determine the training outputs on that cell.
They do not, by themselves, determine a new probe's output, even if every
individual probe sign is specified. For a concrete example use a single
training vector $e_1$, $a,b>0$, and two balanced neurons $w_i=u_i$.
Compare
\[
 u_1=(a,b),\quad u_2=(a,-b)
 \quad\hbox{with}\quad
 u'_1=(7a/5,b/5),\quad u'_2=(a/5,-7b/5).
\]
This is an orthogonal mixing of the two neuron columns. Both training
patterns are active, and in both states the $-e_2$ probe is inactive on
neuron 1 and active on neuron 2. All three training-group moments are
identical, equal to $\operatorname{diag}(2a^2,2b^2)$. Yet
\[
 f(-e_2)=(ab,-b^2),\qquad
 f'(-e_2)=(7ab/25,-49b^2/25).
\]
Thus probe-refined moments (or an equivalent signed probe observable)
are necessary for an OOD interface. This example concerns the sufficiency
of proposed observable data; it is not a counterexample to canonical
initialized reversal.

\paragraph{Gate-event transfers.}
Let $y_i=(u_i,w_i)$ and $K_i=y_iy_i^T$. At an isolated switch of neuron $i$
from training group $g$ to $g'$, the underlying state is continuous, but
its group representation changes by
\[
 K_g^+=K_g^- -K_i,\qquad K_{g'}^+=K_{g'}^-+K_i.
\]
Simultaneous switches sum their individual contributions. The group moment
does not identify $K_i$. Indeed different rank-one decompositions, including
the orthogonal mixing above, have the same $K_g$ and different $K_i$.
An estimate of the form
\begin{equation}\label{eq:invalid-transfer-bound}
 \|K_i-\widehat K_i\|\le C\|K_g-\widehat K_g\|
\end{equation}
therefore fails for every finite universal $C$. A transfer enclosure must
retain the boundary neuron's rank-one moment or other sufficient geometry.
This is an information loss of the unaugmented quotient, not a failure of
the proposed augmented quotient/event strategy.

\paragraph{Reference census.}
The deterministic census of the stored numerical reference has 53 distinct
training patterns at $t=0$, 121 at $t=1$, 140 at $t=2$, 134 at $t=3.4$,
and 134 at $t=4$. A 29-neuron inactive training group persists at these
reference knots. Over 20000 adjacent reference-knot pairs there are 185211
sample--neuron sign differences. These are diagnostic counts, not certified
actual-flow events; multiple crossings between knots could be missed.
At entry a 16-neuron training-pattern group mixes strong-cohort membership.
It must be split by membership to recover the repository $Q_G$ from group
moments. The stepwise CSV also records all group sizes and a nonrigorous
linear-motion near-boundary count. No census number is a proof premise.

\end{document}
